\problemstart
\addcontentsline{toc}{section}{1.2　相対運動と相対速度}

\begin{problemBox}{問題 1.2\quad 相対運動と相対速度\hspace{1.2em}{\small\mdseries 〔基本〕}}
基準となる慣性座標系の原点 O から見た質点1の位置ベクトルを $\boldsymbol{r}_1$ ，質点2の位置ベクトルを $\boldsymbol{r}_2$ とする。質点1から見た質点2の相対位置ベクトルを $\boldsymbol{r}_{12}$ と定義する。以下の問いに答えよ。

\begin{enumerate}
  \item 質点1から見た質点2の相対位置ベクトル $\boldsymbol{r}_{12}$ を， $\boldsymbol{r}_1$ と $\boldsymbol{r}_2$ を用いて表せ。

  \item 質点1の速度ベクトルを $\boldsymbol{v}_1 = \frac{d\boldsymbol{r}_1}{dt}$ ，質点2の速度ベクトルを $\boldsymbol{v}_2 = \frac{d\boldsymbol{r}_2}{dt}$ とする。質点1から見た質点2の相対速度ベクトル $\boldsymbol{v}_{12} = \frac{d\boldsymbol{r}_{12}}{dt}$ を， $\boldsymbol{v}_1$ と $\boldsymbol{v}_2$ を用いて表せ。

  \item 質点1が速度 $\boldsymbol{v}_1 = \begin{pmatrix} u \\ 0 \end{pmatrix}$ で一定の運動を行っており，質点2が速度 $\boldsymbol{v}_2 = \begin{pmatrix} 0 \\ v \end{pmatrix}$ で一定の運動を行っている場合を想定する。初期時刻 $t=0$ における各質点の間隔を $L$ とし，それぞれの初期位置ベクトルを $\boldsymbol{r}_1(0) = \begin{pmatrix} 0 \\ 0 \end{pmatrix}$ ， $\boldsymbol{r}_2(0) = \begin{pmatrix} 0 \\ L \end{pmatrix}$ とする。任意の時刻 $t$ における相対速度ベクトル $\boldsymbol{v}_{12}$ を求めよ。また，2つの質点の間隔が最小となる時刻 $t_{\text{min}}$ を $u, v, L$ を用いて表せ。
\end{enumerate}
\end{problemBox}

\solutionheading
\begin{solutionparts}

\solutionpart{(1)}
ベクトル合成の法則より，原点 O から質点1を経由して質点2へ至る経路を組み立てると， $\boldsymbol{r}_1 + \boldsymbol{r}_{12} = \boldsymbol{r}_2$ という関係が成り立ちます。この等式を相対位置ベクトルについて整理します。
\begin{equation*}
\boldsymbol{r}_{12} = \boldsymbol{r}_2 - \boldsymbol{r}_1
\end{equation*}

\solutionpart{(2)}
(1) で得られた相対位置ベクトルの式を，時間 $t$ で微分します。微分の線形性より，次のようになります。
\begin{align*}
\boldsymbol{v}_{12}
  &= \frac{d\boldsymbol{r}_{12}}{dt}
   = \frac{d}{dt}(\boldsymbol{r}_2-\boldsymbol{r}_1) \\
  &= \frac{d\boldsymbol{r}_2}{dt}
   - \frac{d\boldsymbol{r}_1}{dt}
\end{align*}
それぞれの速度ベクトルの定義を代入します。
\begin{equation*}
\boldsymbol{v}_{12} = \boldsymbol{v}_2 - \boldsymbol{v}_1
\end{equation*}

\solutionpart{(3)}
提示された具体的な速度ベクトルを (2) の結果に代入します。
\begin{align*}
\boldsymbol{v}_{12}
  &= \begin{pmatrix} 0 \\ v \end{pmatrix}
   - \begin{pmatrix} u \\ 0 \end{pmatrix} \\
  &= \begin{pmatrix} -u \\ v \end{pmatrix}
\end{align*}
これが任意の時刻における相対速度ベクトルです。 次に，任意の時刻における相対位置ベクトルを求めます。各質点は等速度運動を行うため，位置ベクトルは次のようになります。
\begin{align*}
\boldsymbol{r}_1(t)
  &= \begin{pmatrix} 0 \\ 0 \end{pmatrix}
   + \begin{pmatrix} u \\ 0 \end{pmatrix}t
   = \begin{pmatrix} ut \\ 0 \end{pmatrix}, \\
\boldsymbol{r}_2(t)
  &= \begin{pmatrix} 0 \\ L \end{pmatrix}
   + \begin{pmatrix} 0 \\ v \end{pmatrix}t
   = \begin{pmatrix} 0 \\ L+vt \end{pmatrix}.
\end{align*}
(1) の定義に従って相対位置ベクトルを求めます。
\begin{align*}
\boldsymbol{r}_{12}(t)
  &= \begin{pmatrix} 0 \\ L+vt \end{pmatrix}
   - \begin{pmatrix} ut \\ 0 \end{pmatrix} \\
  &= \begin{pmatrix} -ut \\ L+vt \end{pmatrix}
\end{align*}
2つの質点の間隔の2乗を $f(t) = |\boldsymbol{r}_{12}(t)|^2$ と置きます。間隔が最小になるとき，その2乗も最小になります。
\begin{align*}
f(t)
  &= (-ut)^2+(L+vt)^2 \\
  &= u^2t^2+L^2+2Lvt+v^2t^2 \\
  &= (u^2+v^2)t^2+2Lvt+L^2
\end{align*}
この $t$ に関する2次関数を平方完成します。
\begin{align*}
f(t)
  &= (u^2+v^2)
     \left(t+\frac{Lv}{u^2+v^2}\right)^2 \\
  &\quad +L^2-\frac{L^2v^2}{u^2+v^2}
\end{align*}
実数の2乗の項が 0 になるときに全体の関数値が最小になります。したがって，間隔が最小となる時刻は次のようになります。
\begin{equation*}
t_{\text{min}} = - \frac{Lv}{u^2 + v^2}
\end{equation*}
実際の物理的な状況において，時刻 $t$ は 0 以上を想定するため， $u, v, L$ が正の値であるならば，この結果は過去の時刻においてすでに間隔が最小であったことを示します。

\end{solutionparts}
